\documentclass[11pt]{article}
\usepackage[margin=0.85in]{geometry}
\usepackage{amsmath,amssymb}
\usepackage{xcolor,listings}
\usepackage{hyperref}
\hypersetup{colorlinks=true,linkcolor=blue}
\lstset{language=Python,basicstyle=\ttfamily\small,columns=fullflexible,
 keepspaces=true,showstringspaces=false,backgroundcolor=\color{black!3},
 frame=single,rulecolor=\color{black!15},aboveskip=7pt,belowskip=7pt}
\setlength{\parindent}{0pt}
\setlength{\parskip}{5pt}
\newcommand{\ket}[1]{\lvert #1\rangle}
\newcommand{\bra}[1]{\langle #1\rvert}
\newcommand{\mtwo}[4]{\begin{pmatrix}#1&#2\\#3&#4\end{pmatrix}}
\newcommand{\vfour}[4]{\begin{pmatrix}#1\\#2\\#3\\#4\end{pmatrix}}
\newcommand{\heading}[1]{\begin{center}\large\textit{#1}\end{center}}
\begin{document}

\heading{Schmidt Decomposition and Truncation}
\textit{Examples 03 and 04}

\subsection*{1. From SVD to Schmidt decomposition}
Begin with a normalized state and factor its coefficient matrix:
\begin{equation}
\ket{\psi}=\sum_{ij}C_{ij}\ket{i}_A\ket{j}_B,
\qquad C=U\Sigma V^\dagger,\qquad s_1\geq s_2\geq\cdots\geq0.
\end{equation}
Substituting gives
\begin{align}
C_{ij}&=\sum_\alpha U_{i\alpha}s_\alpha(V^\dagger)_{\alpha j}
=\sum_\alpha s_\alpha U_{i\alpha}V_{j\alpha}^*,\\
\ket{\psi}&=\sum_\alpha s_\alpha
\underbrace{\left(\sum_i U_{i\alpha}\ket{i}_A\right)}_{\ket{a_\alpha}}
\otimes\underbrace{\left(\sum_j V_{j\alpha}^*\ket{j}_B\right)}_{\ket{b_\alpha}}.
\end{align}
Since $U^\dagger U=V^\dagger V=I$ on the retained columns,
\begin{equation}
\langle a_\alpha|a_\beta\rangle
=\langle b_\alpha|b_\beta\rangle=\delta_{\alpha\beta},
\qquad \ket{\psi}=\sum_{\alpha=1}^r s_\alpha\ket{a_\alpha}\ket{b_\alpha},
\qquad\sum_\alpha s_\alpha^2=1.
\end{equation}
The Schmidt rank is $r=\operatorname{rank}C$. A product state has $r=1$.

In \texttt{U, s, Vh = np.linalg.svd(C, full\_matrices=False)}, the ket coefficients are
\begin{equation}
\ket{a_\alpha}\ \longleftrightarrow\ \texttt{U[:, alpha]},
\qquad \ket{b_\alpha}\ \longleftrightarrow\ \texttt{Vh[alpha, :]}.
\end{equation}
The row of \texttt{Vh} already contains the required conjugation.

\subsection*{2. Reduced density matrices and entropy}
Tracing over one subsystem gives
\begin{align}
(\rho_A)_{ii'}&=\sum_j C_{ij}C^*_{i'j},&\rho_A&=CC^\dagger,\\
(\rho_B)_{jj'}&=\sum_i C_{ij}C^*_{ij'},&\rho_B&=C^{\mathsf T}C^*.
\end{align}
These are \texttt{C @ C.conj().T} and \texttt{C.T @ C.conj()}. Using the Schmidt form,
\begin{align}
\rho_A&=\sum_{\alpha\beta}s_\alpha s_\beta\ket{a_\alpha}\bra{a_\beta}
\langle b_\beta|b_\alpha\rangle
=\sum_\alpha s_\alpha^2\ket{a_\alpha}\bra{a_\alpha},\\
\rho_B&=\sum_\alpha s_\alpha^2\ket{b_\alpha}\bra{b_\alpha}.
\end{align}
Both have nonzero eigenvalues $p_\alpha=s_\alpha^2$, so
\begin{equation}
S_A=S_B=-\sum_\alpha p_\alpha\log_2p_\alpha.
\end{equation}
The product state has $p=(1,0)$ and $S=0$; the Bell state has $p=(1/2,1/2)$ and $S=1$ bit. Terms with $p=0$ contribute zero; the code excludes numerical zeros.

\newpage
\subsection*{3. Truncation, discarded weight, and fidelity}
Keep the largest $\chi$ Schmidt coefficients:
\begin{equation}
\ket{\widetilde\psi_\chi}=\sum_{\alpha=1}^\chi s_\alpha\ket{a_\alpha}\ket{b_\alpha},
\qquad w=\sum_{\alpha>\chi}s_\alpha^2.
\end{equation}
Orthogonality gives
\begin{align}
\|\psi-\widetilde\psi_\chi\|^2
&=\left\|\sum_{\alpha>\chi}s_\alpha\ket{a_\alpha}\ket{b_\alpha}\right\|^2
=\sum_{\alpha>\chi}s_\alpha^2=w,\\
\langle\widetilde\psi_\chi|\widetilde\psi_\chi\rangle
&=\sum_{\alpha\leq\chi}s_\alpha^2=1-w.
\end{align}
The truncated SVD minimizes this error among coefficient matrices of rank at most $\chi$.

Normalize the truncated state:
\begin{equation}
\ket{\phi_\chi}=\frac{\ket{\widetilde\psi_\chi}}{\sqrt{1-w}}.
\end{equation}
Then
\begin{align}
\langle\psi|\phi_\chi\rangle
&=\frac{\sum_{\alpha\leq\chi}s_\alpha^2}{\sqrt{1-w}}
=\sqrt{1-w},\\
F&=|\langle\psi|\phi_\chi\rangle|^2=1-w,\\
\|\psi-\phi_\chi\|^2&=2-2\operatorname{Re}\langle\psi|\phi_\chi\rangle
=2(1-\sqrt{1-w}).
\end{align}
Thus discarded weight measures the error before normalization and the loss of fidelity after normalization.

\subsection*{4. The two-spin example}
For $0\leq\epsilon\leq1/2$, take
\begin{equation}
\ket{\psi_\epsilon}=\sqrt{1-\epsilon}\ket{\uparrow\uparrow}
+\sqrt\epsilon\ket{\downarrow\downarrow}.
\end{equation}
Keeping only the first Schmidt term gives
\begin{align}
\ket{\widetilde\psi}&=\sqrt{1-\epsilon}\ket{\uparrow\uparrow}
\quad\xrightarrow{\text{normalize}}\quad\ket{\phi}=\ket{\uparrow\uparrow},\\
w&=\epsilon,\qquad F=1-\epsilon,\\
S&=-(1-\epsilon)\log_2(1-\epsilon)-\epsilon\log_2\epsilon.
\end{align}
For $\epsilon=0.05$, the squared error before normalization is $0.05$, and the fidelity afterward is $0.95$.

\newpage
\subsection*{Observable errors}
Use
\begin{equation}
Z=\mtwo{1}{0}{0}{-1},\qquad X=\mtwo{0}{1}{1}{0}.
\end{equation}
Acting on the original state,
\begin{align}
(Z\otimes I)\ket{\psi_\epsilon}
&=\sqrt{1-\epsilon}\ket{\uparrow\uparrow}-\sqrt\epsilon\ket{\downarrow\downarrow},\\
(Z\otimes Z)\ket{\psi_\epsilon}&=\ket{\psi_\epsilon},\\
(X\otimes X)\ket{\psi_\epsilon}
&=\sqrt{1-\epsilon}\ket{\downarrow\downarrow}+\sqrt\epsilon\ket{\uparrow\uparrow}.
\end{align}
Taking the inner products gives
\begin{align}
\langle Z\otimes I\rangle_{\psi_\epsilon}&=(1-\epsilon)-\epsilon=1-2\epsilon,
&\langle Z\otimes I\rangle_\phi&=1,\\
\langle Z\otimes Z\rangle_{\psi_\epsilon}&=1,
&\langle Z\otimes Z\rangle_\phi&=1,\\
\langle X\otimes X\rangle_{\psi_\epsilon}&=2\sqrt{\epsilon(1-\epsilon)},
&\langle X\otimes X\rangle_\phi&=0.
\end{align}
Hence
\begin{equation}
\Delta_{Z\otimes I}=2\epsilon,\qquad
\Delta_{Z\otimes Z}=0,\qquad
\Delta_{X\otimes X}=2\sqrt{\epsilon(1-\epsilon)}.
\end{equation}
The same truncation leaves one observable unchanged and alters the others. Fidelity describes the whole state; an individual observable can be much more accurate.

For a general Hermitian $O$, use the standard trace-norm identities
\begin{equation}
|\operatorname{Tr}(O D)|\leq\|O\|_{\rm op}\|D\|_1,
\qquad
\bigl\|\ket{\psi}\bra{\psi}-\ket{\phi_\chi}\bra{\phi_\chi}\bigr\|_1
=2\sqrt{1-F}.
\end{equation}
Substituting gives
\begin{align}
|\langle O\rangle_\psi-\langle O\rangle_{\phi_\chi}|
&=\left|\operatorname{Tr}\!\left[O\left(\ket{\psi}\bra{\psi}
-\ket{\phi_\chi}\bra{\phi_\chi}\right)\right]\right|\\
&\leq2\|O\|_{\rm op}\sqrt{1-F}=2\|O\|_{\rm op}\sqrt w.
\end{align}
Here $\|O\|_{\rm op}$ is the largest absolute eigenvalue, and $\|D\|_1$ is the sum of the singular values of $D$.

\end{document}